\documentclass[10pt,twocolumn]{article}

% ============================================================
% PAGE AND TYPOGRAPHY
% ============================================================
\usepackage[a4paper,margin=1.7cm]{geometry}
\usepackage[T1]{fontenc}
\usepackage[utf8]{inputenc}
\usepackage{lmodern}
\usepackage{microtype}
\usepackage{setspace}
\usepackage{tabularx}
\usepackage{ragged2e}
% \usepackage{cuted} % REMOVED: strip/cuted rewires the two-column output routine and can
% produce empty-column artifacts with floats in this manuscript.
% Do not auto-insert \FloatBarrier at every \section in two-column mode.
\usepackage{placeins}

% ============================================================
% MATHEMATICS, TABLES AND FIGURES
% ============================================================
\usepackage{amsmath,amssymb}
\usepackage{graphicx}
\usepackage{booktabs}
\usepackage{multirow}
\usepackage{array}
\usepackage{float}
\usepackage{caption}
\usepackage{subcaption}
\renewcommand{\topfraction}{0.9}
\renewcommand{\bottomfraction}{0.8}
\renewcommand{\textfraction}{0.07}
\renewcommand{\floatpagefraction}{0.8}

% ============================================================
% REFERENCES -- APA 7 + CLICKABLE CITATIONS AND DOI LINKS
% ============================================================
\usepackage[american]{babel}
\usepackage{csquotes}

\usepackage[
    backend=biber,
    style=apa,
    sorting=nyt,
    maxcitenames=2,
    maxbibnames=20,
    doi=true,
    url=true,
    isbn=false
]{biblatex}
\usepackage{xcolor}

\addbibresource{references.bib}

% ============================================================
% HYPERLINKS
% ============================================================
\usepackage[
    hidelinks,
    bookmarks=true,
    bookmarksnumbered=true,
    colorlinks=true,
    linkcolor=black,
    citecolor=blue,
    urlcolor=blue
]{hyperref}

% ============================================================
% BIBLIOGRAPHY FORMATTING
% ============================================================
\setlength{\bibitemsep}{0.55\baselineskip}
\setlength{\bibhang}{0.25in}

% ============================================================
% COLUMN AND SECTION FORMATTING
% ============================================================
\usepackage{titlesec}

\titleformat{\section}
  {\normalfont\bfseries\large}
  {\thesection.}
  {0.5em}
  {}

\titleformat{\subsection}
  {\normalfont\bfseries\normalsize}
  {\thesubsection}
  {0.5em}
  {}

\titleformat{\subsubsection}
  {\normalfont\itshape\normalsize}
  {\thesubsubsection}
  {0.5em}
  {}

\titlespacing*{\section}{0pt}{8pt}{4pt}
\titlespacing*{\subsection}{0pt}{6pt}{3pt}
\titlespacing*{\subsubsection}{0pt}{4pt}{2pt}

% ============================================================
% DOCUMENT
% ============================================================

\begin{document}

% ============================================================
% FIRST PAGE HEADER + ABSTRACT
% ============================================================

% Full-width first-page block implemented without cuted/strip
% (must avoid literal ']' tokens inside the argument).
\twocolumn[

% ------------------------------------------------------------
% TITLE
% ------------------------------------------------------------

\noindent
{\fontsize{17}{20}\selectfont\bfseries
Biologically Informed Machine Learning Framework for Malaria Incidence
Prediction Using Vectorial-Capacity-Derived Information
\par}

\vspace{0.8em}

% ------------------------------------------------------------
% AUTHORS
% ------------------------------------------------------------

\noindent
{\large\bfseries
Kodjo Gato Didier EWLI$^{\hyperlink{aff:a}{\textcolor{blue}{a}},*}$,
Mwangi Karanja$^{\hyperlink{aff:b}{\textcolor{blue}{b}}}$,
Komi Mensah Agboka$^{\hyperlink{aff:c}{\textcolor{blue}{c}}}$,
Antoine Barreaux$^{\hyperlink{aff:c}{\textcolor{blue}{c}}}$,
Folorunso Sakinat Oluwabukonla$^{\hyperlink{aff:d}{\textcolor{blue}{d}}}$,
Selina Akouwa WUIGA$^{\hyperlink{aff:e}{\textcolor{blue}{e}}}$
\par}

\vspace{0.5em}

% ------------------------------------------------------------
% AFFILIATIONS
% ------------------------------------------------------------

\noindent
{\footnotesize
\hypertarget{aff:a}{$^{a}$} Pan African University Institute for Basic Sciences,
Technology and Innovation (PAUSTI), Nairobi, Kenya\\
\hypertarget{aff:b}{$^{b}$} Jomo Kenyatta University of Agriculture and Technology
(JKUAT), Nairobi, Kenya\\
\hypertarget{aff:c}{$^{c}$} International Centre of Insect Physiology and Ecology
(icipe), P.O. Box 30772-00100, Nairobi, Kenya\\
\hypertarget{aff:d}{$^{d}$} Olabisi Onabanjo University, Ago-Iwoye, Nigeria\\
\hypertarget{aff:e}{$^{e}$} Dibrugarh University, 786004, Assam, India\\\vspace{2pt}
$^{*}$ Corresponding author: Kodjo Gato Didier EWLI
\par}

\vspace{0.6em}

\hrule

\vspace{0.7em}

% ------------------------------------------------------------
% ARTICLE INFO + ABSTRACT
% ------------------------------------------------------------

\noindent
\begin{minipage}{0.25\textwidth}
\vspace{0pt}
\footnotesize

\textbf{A R T I C L E \quad I N F O}

\vspace{0.25em}
\hrule
\vspace{0.45em}

\textit{Keywords:}

\vspace{0.25em}

Malaria incidence prediction\\\vspace{2pt}
Biologically informed machine learning\\\vspace{2pt}
Vectorial capacity\\\vspace{2pt}
Random Forest\\\vspace{2pt}
Gradient Boosting\\\vspace{2pt}
Remote sensing\\\vspace{2pt}
Sub-Saharan Africa

\end{minipage}
\hfill
\begin{minipage}{0.70\textwidth}
\vspace{0pt}
\footnotesize

\textbf{A B S T R A C T}

\vspace{0.25em}
\hrule
\vspace{0.45em}

Malaria transmission is shaped by complex interactions among
epidemiological, environmental, and vector-related processes, yet
machine-learning prediction models often represent these relationships
primarily through statistical associations. This study evaluated
whether vectorial-capacity-derived information provides incremental
predictive value beyond historical malaria incidence and environmental
predictors within a biologically informed machine-learning framework.
Annual country-level data from 32 Sub-Saharan African countries were
constructed from malaria incidence, land-surface temperature, rainfall,
vegetation, surface-moisture, and population-density information. The
target-complete dataset contained 159 country-year observations, of
which 127 were eligible for the primary one-lag analysis. Vectorial
capacity (VC) was represented through an environmental parameterization
of the classical formulation, and VC-derived features were evaluated
using Random Forest and Gradient Boosting regression models within a
structured feature-ablation design. Performance was assessed using
repeated country-grouped validation and a separate 2020 temporal
holdout.

Historical incidence remained the dominant predictive signal. Under
country-grouped validation, several machine-learning configurations
reduced RMSE and increased \(R^2\) relative to incidence persistence,
although persistence retained the lowest mean MAE. In the 2020
temporal evaluation, the one-lag history-plus-VC Random Forest model
achieved the lowest RMSE among the evaluated configurations, with an
MAE of 30.5827, RMSE of 50.4361, and \(R^2\) of 0.7075. Relative to
persistence, this represented reductions of approximately 5.2\% in
MAE and 7.6\% in RMSE. More importantly, relative to the corresponding
history-only Random Forest model, incorporating VC-derived information
reduced temporal RMSE by approximately 2.33 incidence points and MAE
by 0.99 incidence points, demonstrating incremental predictive value
beyond historical incidence alone in this configuration.

The predictive contribution of VC-derived information was
context-dependent and varied across feature representations,
countries, and evaluation metrics. Its clearest incremental value was
observed when combined with historical incidence in the one-lag Random
Forest model. These findings provide a proof-of-concept that
environmentally parameterized VC can serve as a compact, biologically
structured representation that complements historical incidence within
malaria prediction, while highlighting the need for further validation
using larger, higher-resolution datasets and direct entomological
observations.

\end{minipage}

\vspace{0.7em}

\hrule

]


% ============================================================
% NORMAL TWO-COLUMN BODY STARTS HERE
% ============================================================

\section{Introduction}

\subsection{Background}

Malaria remains a major vector-borne disease whose transmission is
shaped by interactions among parasite biology, mosquito ecology,
environmental conditions, human populations, and public-health
interventions. Climatic and ecological factors such as temperature,
rainfall, vegetation, and land-surface conditions influence mosquito
survival and breeding, parasite development, and environmental
suitability for transmission
\parencite{LindsayBirley1996,MafweleLee2022,Odhiambo2020}.

Temperature affects mosquito development, survival, biting behaviour,
and parasite development within the vector
\parencite{Villena2022,Suh2024}, while rainfall and vegetation provide
information on breeding conditions and habitat suitability.
Consequently, environmental variability contributes to spatial and
temporal heterogeneity in malaria transmission
\parencite{Ferraguti2023,Rakotoarinia2022,Dahlin2024,
deSouzaWeaver2024}. Remote-sensing observations provide a practical
means of characterizing these environmental conditions across large
geographical areas and have increasingly been used in malaria-risk
assessment and prediction \parencite{Cuellar2025}.

These relationships have motivated the use of machine learning for
malaria prediction because such models can capture nonlinear
relationships and interactions among epidemiological and environmental
variables \parencite{Masinde2020,Odu2021,Khan2024,
RahmanShiddik2025}. However, predictive associations learned from
observational data do not necessarily represent the biological
processes underlying transmission. This limitation has encouraged
biologically informed and knowledge-guided machine learning, in which
domain knowledge is incorporated into feature construction or model
development \parencite{Wang2020,Yin2019,Jhutty2022,Karpatne2024}.

One biologically meaningful representation of malaria transmission
potential is vectorial capacity (VC), which characterizes the potential
contribution of a mosquito population to pathogen transmission through
vector-related biological processes. Its foundations originate from
classical malaria-transmission theory \parencite{Ross1905,Ross1911a},
with subsequent work extending its application to environmental
monitoring and malaria surveillance \parencite{Ceccato2012}. VC is
therefore distinct from observed malaria incidence and can provide
complementary biological information on transmission potential.

This study investigates whether VC-derived information, constructed
through an environmental parameterization of the classical VC
framework, provides incremental predictive information when combined
with historical malaria incidence and environmental predictors. The
objective is not to propose a new VC formulation or replace observed
incidence, but to evaluate VC as a biologically informed representation
of transmission potential within a machine-learning framework.


\subsection{Research Gap}

Although machine learning has demonstrated potential for malaria
prediction, the incremental value of biologically informed transmission
information remains insufficiently evaluated when strong historical
incidence and environmental predictors are already available
\parencite{Odu2021,Okiring2021,MafweleLee2022,Kaur2022,Mbunge2023}.
This distinction is important because malaria incidence exhibits
substantial temporal persistence. Consequently, the predictive value
of VC-derived information should be assessed beyond the information
already contained in recent incidence and conventional environmental
predictors.

A related challenge concerns interpretation. Incorporating
biologically related environmental information does not demonstrate
that a machine-learning model has learned a biological mechanism, and
feature importance should not be interpreted as evidence of causality
\parencite{FanLiuSun2023,Qiu2022,AlMeslamani2024}. The contribution of
VC should therefore be evaluated through controlled predictive
comparisons rather than feature importance alone.

Accordingly, the central research question is:

\begin{quote}
\textit{Does incorporating vectorial-capacity-derived information into
a machine-learning framework provide incremental predictive value
beyond historical malaria incidence and environmental predictors?}
\end{quote}

This question is examined using country-level observations from
sub-Saharan Africa under both country-grouped validation and a
temporal evaluation using 2020 as the held-out observation period.


\subsection{Research Objective}

The overall objective of this study is to evaluate whether the
integration of vectorial-capacity-derived information with historical
incidence and environmental predictors can improve machine-learning
prediction of malaria incidence.

The specific objectives are to:

\begin{enumerate}
    \item construct a biologically informed representation of malaria
    transmission potential using environmental and demographic
    information related to vectorial capacity;

    \item integrate vectorial-capacity-derived features with historical
    incidence and environmental predictors within a machine-learning
    framework;

    \item evaluate Random Forest and Gradient Boosting models under
    country-grouped and temporal evaluation settings;

    \item compare the models against an incidence-persistence baseline
    based on the last available incidence observation;

    \item quantify predictive performance using mean absolute error
    (MAE), root mean squared error (RMSE), and coefficient of
    determination ($R^2$); and

    \item examine the predictive contribution of vectorial-capacity-derived
    features without interpreting feature importance as evidence of
    causality.
\end{enumerate}


\subsection{Main Contribution}

The principal contribution of this study is empirical rather than
theoretical. Rather than proposing a new vectorial-capacity equation
or mechanistic malaria model, the study evaluates whether an
environmentally parameterized VC representation provides incremental
predictive information beyond historical incidence and conventional
environmental predictors.

The framework integrates three complementary sources of information:
historical incidence representing temporal persistence; environmental
and demographic predictors representing ecological and population
conditions; and VC-derived features providing a biologically informed
representation of transmission potential. Their contributions are
evaluated through controlled feature-ablation experiments under
country-grouped and temporal evaluation.

The study therefore provides a direct assessment of whether biological
structure contributes additional predictive information beyond strong
epidemiological and environmental signals. Improvements are interpreted
as evidence of predictive utility rather than proof of causal
biological mechanisms.


% ============================================================
% SECTION 2 WILL FOLLOW
% ============================================================

% \section{Related Work}
% This section will be added in the next step.
% ============================================================
% 2. RELATED WORK
% ============================================================


\section{Related Work}

\subsection{Environmental and Climate-Based Malaria Prediction}

Environmental and climatic conditions influence malaria transmission
through their effects on the vector--parasite--host system. Temperature
affects mosquito development, survival, and parasite development,
whereas rainfall, vegetation, and humidity influence breeding
conditions, habitat suitability, and vector survival
\parencite{LindsayBirley1996,MafweleLee2022,Dahlin2024}. These
relationships are frequently nonlinear, particularly for temperature,
because increasing temperature does not continuously increase
transmission suitability across all ranges
\parencite{Villena2022,Suh2024}.

Remote sensing provides a practical means of characterizing these
environmental conditions across large geographical areas. Vegetation,
land-surface, and related environmental observations have been
associated with spatial variation in vector habitats and malaria risk
\parencite{Ferraguti2023,Sheela2017,Rakotoarinia2022}. Recent studies
have further demonstrated the value of remotely sensed environmental
information for malaria-risk assessment and prediction
\parencite{Cuellar2025}.

Environmental effects may also precede changes in observed incidence
because mosquito development, parasite development, and disease
reporting operate over different time scales. Lagged environmental
relationships have therefore been investigated in malaria modelling
\parencite{Okiring2021}. Nevertheless, environmental predictors alone
cannot represent the complete epidemiological process, which is also
influenced by population characteristics, interventions, healthcare
access, surveillance, and vector ecology.


\subsection{Machine Learning for Malaria Incidence and Outbreak Prediction}

Machine learning has increasingly been applied to malaria prediction
because of its ability to represent nonlinear relationships and
interactions among climatic, environmental, demographic, and
epidemiological predictors. Studies in African settings have
demonstrated its application to malaria-incidence prediction using
climate and epidemiological information
\parencite{Masinde2020,Odu2021}, while more recent work has extended
these approaches to outbreak prediction and large-scale malaria
analysis \parencite{Khan2024,RahmanShiddik2025}. Reviews of
artificial-intelligence applications to vector-borne diseases also
identify data quality, interpretability, and generalizability as
important challenges \parencite{Kaur2022}.

A particularly important consideration is temporal persistence in
malaria incidence. Historical incidence may contain substantial
predictive information because epidemiological, environmental, and
health-system characteristics can persist over time. Consequently,
claims about the contribution of additional environmental or biological
information should be evaluated against a strong historical-incidence
reference. The present study therefore uses an incidence-persistence
baseline based on the \emph{last-observed incidence}.


\subsection{Vectorial Capacity as a Biological Representation of
Transmission Potential}

Vectorial capacity provides an established epidemiological framework
for relating mosquito characteristics and transmission processes to
the potential for pathogen transmission. Its theoretical foundations
originate from classical mathematical studies of malaria transmission
\parencite{Ross1905,Ross1911a}, with subsequent work demonstrating how
VC-related quantities can be combined with environmental information
for malaria surveillance and transmission monitoring
\parencite{Ceccato2012}.

The biological basis of VC is supported by evidence that environmental
conditions affect several components of mosquito-borne parasite
transmission. Temperature influences vector survival, transmission
suitability, and \textit{Plasmodium} development within mosquitoes
\parencite{Villena2022,Suh2024}, while environmental and host
conditions can further modify transmission potential
\parencite{Dahlin2024}.

In this study, VC is used as a biologically informed, environmentally
parameterized representation of transmission potential rather than as
a direct entomological measurement or replacement for observed
incidence. VC-derived features are incorporated alongside historical
incidence and environmental predictors to test whether this biological
representation provides incremental predictive information.


\subsection{Biologically Informed and Mechanism-Guided Machine
Learning}

Knowledge-guided and mechanistically informed machine learning seeks
to integrate scientific knowledge with data-driven prediction
\parencite{Wang2020,Yin2019,Karpatne2024}. In epidemiology, hybrid
approaches have combined machine-learning flexibility with
mechanistic or domain-specific representations
\parencite{Jhutty2022,Han2024,Millevoi2024}. Such approaches are
particularly relevant when observational data are limited,
heterogeneous, or insufficient to recover the underlying transmission
process directly.

However, biological plausibility does not guarantee additional
predictive value. A biologically meaningful feature may contribute
little once correlated environmental variables and historical
incidence are included. Conversely, a biologically structured
representation may summarize transmission information not efficiently
captured by individual environmental predictors. Recent work therefore
emphasizes the importance of empirical validation when integrating
mechanistic and data-driven approaches \parencite{Ye2025}.

Interpretability must also be distinguished from causality. Feature
importance describes how a predictor is used by a fitted model but does
not establish that the predictor causes the observed outcome,
particularly when environmental and biologically derived variables are
correlated \parencite{FanLiuSun2023,Qiu2022,AlMeslamani2024}. This
motivates the use of controlled feature-ablation comparisons as the
primary assessment of VC contribution in the present study.


\subsection{Data Quality and Epidemiological Heterogeneity}

Malaria surveillance data may be affected by underreporting, delayed
reporting, diagnostic differences, and variation in surveillance
coverage. Research on infectious-disease modelling has demonstrated
the potential effects of underreporting and reporting delays on
observed disease data and predictive performance
\parencite{Arima2024,Beesley2022}. Consequently, observed malaria
incidence represents an epidemiological measurement that may not
perfectly reflect underlying transmission intensity.

Malaria transmission is also spatially heterogeneous because
environmental suitability, vector populations, human characteristics,
interventions, and healthcare conditions vary geographically
\parencite{Odhiambo2020,Wu2023}. This heterogeneity motivates
country-grouped evaluation and country-level examination of prediction
errors, since improvements in aggregate performance may not occur
uniformly across countries.


\subsection{Research Gap and Positioning of the Present Study}

Previous research establishes that environmental and remote-sensing
variables contain information relevant to malaria transmission
\parencite{MafweleLee2022,Okiring2021,Cuellar2025}, that machine
learning can exploit nonlinear environmental--epidemiological
relationships \parencite{Odu2021,Masinde2020,Khan2024}, and that
knowledge-guided approaches can incorporate biological understanding
into data-driven prediction \parencite{Jhutty2022,Karpatne2024,Ye2025}.

What remains insufficiently established is whether a
vectorial-capacity-derived representation provides additional
predictive information once historical incidence and conventional
environmental predictors are already available. The present study
addresses this question through controlled feature-ablation experiments
using Random Forest and Gradient Boosting models under country-grouped
and temporal evaluation. The aim is not to establish a new VC equation
or causal mechanism, but to determine whether biologically informed
VC-derived information provides measurable incremental predictive
value.


%============================================
%Section 3: Methods ans Materials
%============================================

\section{Materials and Methods}

\subsection{Study Design and Analytical Dataset}

This study evaluated whether incorporating biologically informed
vectorial-capacity-derived information into a machine-learning
framework improves malaria-incidence prediction beyond historical
incidence and conventional environmental predictors. Random Forest
(RF) and Gradient Boosting (GB) regression models were evaluated using
historical incidence, environmental and demographic predictors, and
VC-derived features.

The dataset comprised 32 Sub-Saharan African countries with available
observations for 2015, 2016, 2017, 2018, and 2020. Environmental data
were processed using Google Earth Engine \parencite{Gorelick2017}.
Landsat 8 provided land surface temperature (LST), Normalized
Difference Vegetation Index (NDVI), and Normalized Difference Water
Index (NDWI); rainfall was obtained from CHIRPS
\parencite{Funk2015}; population density from WorldPop
\parencite{Tatem2017}; and malaria incidence from the WHO Global
Health Observatory \parencite{WHO_GHO_Malaria}. Malaria incidence,
expressed as cases per 1,000 population at risk, was the prediction
target.

Environmental observations were initially extracted monthly and
aggregated to annual country-level summaries representing average
conditions and within-year variability. Because simultaneous
processing of all countries exceeded practical computational limits,
environmental extraction was performed in batches and subsequently
combined into a single analytical dataset.

The initial dataset contained 160 country-year observations. The year
2019 was absent, and malaria incidence was missing for Mali in 2020,
leaving 159 observations with complete target information. Historical
incidence was constructed chronologically within each country. The
first lag represented the \emph{last-observed incidence}, rather than
necessarily the previous calendar year; consequently, the historical
incidence associated with a 2020 target was the available 2018
observation.

Two analytical designs were evaluated. The one-lag design required
only the last-observed incidence and contained 127 eligible target
observations: 32 each from 2016, 2017, and 2018, and 31 from 2020.
The two-lag design additionally required the second most recent
incidence observation and contained 95 targets: 32 each from 2017 and
2018 and 31 from 2020. Comparing these designs allowed the predictive
value of additional historical information to be assessed against the
loss of 32 eligible observations.

Models were evaluated under two complementary settings.
Country-grouped validation assessed generalization to held-out
countries while allowing their previously observed incidence to remain
available as historical information. A separate temporal evaluation
used 2020 as the held-out period, with models developed from eligible
pre-2020 observations and evaluated on the 31 countries with observed
2020 incidence. Because contemporaneous 2020 environmental predictors
were included, this experiment is interpreted as a temporal holdout
rather than a strictly prospective forecast.

Figure~\ref{fig:biml_framework} summarizes the complete analytical
framework, from environmental and epidemiological information through
VC construction, feature ablation, machine-learning modelling, and
grouped and temporal evaluation.

\begin{figure*}[!t]
    \centering
    \includegraphics[width=\textwidth]{Malaria_Transmission_Modelling_Workflow.png}
    \caption{Biologically informed machine-learning framework for
    malaria-incidence prediction. Environmental, biological, demographic,
    and epidemiological information is integrated through an
    environmentally parameterized vectorial-capacity representation.
    VC-derived features are combined with historical incidence and
    environmental predictors through the M1--M6 ablation design and
    evaluated using Random Forest and Gradient Boosting under
    country-grouped and 2020 temporal evaluation.}
    \label{fig:biml_framework}
\end{figure*}


\subsection{Environmental and Demographic Predictors}

Environmental and demographic predictors represented the principal
ecological conditions associated with malaria transmission, including
temperature, vegetation, surface moisture, rainfall, wet conditions,
and population density. Temperature was characterized using annual
summaries of monthly LST, while NDVI and NDWI represented vegetation
and surface-moisture conditions, respectively. Rainfall variables
characterized hydrological conditions relevant to mosquito breeding,
and population density represented the human population context
\parencite{MafweleLee2022,Okiring2021,Cuellar2025}.

Two environmental representations were evaluated. The \emph{full
environmental representation} retained the annual environmental and
demographic summaries together with the nonlinear features described
below. The \emph{compact environmental representation} was designed to
retain the principal environmental domains while reducing redundancy
and initially included \texttt{Temp\_mean}, \texttt{Temp\_std},
\texttt{NDVI\_mean}, \texttt{NDVI\_std}, \texttt{NDWI\_mean},
\texttt{Rainfall\_total}, \texttt{Rainfall\_std},
\texttt{Wet\_months}, and \texttt{Population\_log}.

Redundancy within the compact representation was assessed using
pairwise Pearson correlations estimated from the training data. For
predictor pairs with \(\lvert r\rvert \geq 0.90\), one variable was
removed according to a predefined feature-priority order. Feature
selection was repeated within each training partition and the retained
set was subsequently applied to the corresponding evaluation data,
preventing held-out observations from influencing predictor selection.


\subsection{Nonlinear Environmental Features}

To represent potential nonlinearities and interactions in
environmental relationships with malaria transmission, six engineered
features were included in the full environmental representation:
\texttt{Rainfall\_log}, \texttt{Population\_log},
\texttt{Temp\_sq}, \texttt{Rainfall\_sq},
\texttt{Temp\_Rain}, and \texttt{NDVI\_NDWI}. They were defined as

\begin{align}
\text{Rainfall}_{\log}
&= \log(1+\text{Rainfall}_{total}), \\
\text{Population}_{\log}
&= \log(1+\text{Population density}), \\
\text{Temp}_{sq}
&= (\text{Temp}_{mean})^2, \\
\text{Rainfall}_{sq}
&= (\text{Rainfall}_{total})^2, \\
\text{Temp\_Rain}
&= \text{Temp}_{mean}\times\text{Rainfall}_{total}, \\
\text{NDVI\_NDWI}
&= \text{NDVI}_{mean}\times\text{NDWI}_{mean}.
\end{align}

The logarithmic transformations reduced differences in the numerical
scale of rainfall and population density, while the squared and
interaction terms represented nonlinear and joint environmental
conditions. Their inclusion was motivated by evidence that malaria
transmission processes, particularly temperature-dependent mosquito
and parasite processes, can respond nonlinearly to environmental
conditions \parencite{Villena2022,Suh2024,Dahlin2024}. These engineered
features were restricted to the full representation; the compact
representation was retained as a smaller, less redundant alternative.



\subsection{Vectorial Capacity Formulation}

Vectorial capacity (VC) was incorporated as a biologically informed
representation of malaria transmission potential. The formulation
follows the classical VC framework, in which transmission potential
depends on mosquito density relative to humans, biting frequency,
mosquito survival, and parasite development within the vector
\parencite{Ross1905,Ross1907,LindsayBirley1996}:

\begin{equation}
C =
\frac{m a^2 p^n}{-\ln(p)},
\label{eq:vc}
\end{equation}

where \(C\) denotes vectorial capacity, \(m\) mosquito density relative
to the human population, \(a\) mosquito biting rate, \(p\) daily
mosquito survival probability, and \(n\) the extrinsic incubation
period of the malaria parasite.

Because direct country-year entomological measurements of these
components were unavailable, \(m\), \(a\), \(p\), and \(n\) were
environmentally parameterized using temperature, rainfall, vegetation,
and population information. The resulting VC representation was
evaluated as biologically structured information within the
machine-learning feature space.


\subsubsection{Environmental Scaling for VC Construction}

Four intermediate quantities represented thermal, rainfall, vegetation,
and population conditions. Thermal suitability was defined as

\begin{equation}
s_T =
\exp\left[
-\frac{(T-25)^2}{2(5)^2}
\right],
\label{eq:thermal_suitability}
\end{equation}

where \(T\) is annual mean land surface temperature. Rainfall
suitability was defined as

\begin{equation}
s_R =
1-\exp\left(
-\frac{R}{R_{\mathrm{med}}^{\mathrm{train}}}
\right),
\label{eq:rain_suitability}
\end{equation}

where \(R\) is annual total rainfall and
\(R_{\mathrm{med}}^{\mathrm{train}}\) is the training-partition median.

Vegetation suitability was obtained by min--max scaling annual mean
NDVI:

\begin{equation}
s_V =
\frac{
V-V_{\min}^{\mathrm{train}}
}{
V_{\max}^{\mathrm{train}}-V_{\min}^{\mathrm{train}}
},
\label{eq:vegetation_suitability}
\end{equation}

where \(V\) denotes annual mean NDVI. Population density was
log-transformed and similarly scaled:

\begin{equation}
P_{\log}=\log(1+P),
\end{equation}

\begin{equation}
s_P =
\frac{
P_{\log}-P_{\log,\min}^{\mathrm{train}}
}{
P_{\log,\max}^{\mathrm{train}}
-
P_{\log,\min}^{\mathrm{train}}
}.
\label{eq:population_suitability}
\end{equation}

All empirical scaling quantities were estimated from the corresponding
training partition. Values outside the training range were bounded to
\([0,1]\).


\subsubsection{Environmental Parameterization of VC Components}

Thermal suitability was used to parameterize biting rate, mosquito
survival, and extrinsic incubation period as

\begin{equation}
a = 0.05 + 0.03s_T,
\label{eq:biting_rate}
\end{equation}

\begin{equation}
p = 0.70 + 0.25s_T,
\label{eq:survival}
\end{equation}

and

\begin{equation}
n = 20 - 8s_T.
\label{eq:eip}
\end{equation}

The mosquito-density component incorporated rainfall, vegetation, and
population conditions:

\begin{equation}
m =
0.25 +
0.75s_Rs_V
\left(
0.5+0.5s_P
\right).
\label{eq:mosquito_density}
\end{equation}

The resulting \(m\), \(a\), \(p\), and \(n\) values were substituted
into Equation~\ref{eq:vc} to calculate \(C\) for each country-year
observation.

These numerical mappings constitute an \emph{operational environmental
parameterization} of the classical VC components for the predictive
experiment. They should not be interpreted as direct country-specific
entomological measurements of mosquito density, biting rate, survival,
or parasite incubation period.


\subsubsection{Normalization and Temporal VC Features}

Calculated VC values were normalized within each training partition as

\begin{equation}
VC_{\mathrm{norm}}
=
\frac{
C-C_{\min}^{\mathrm{train}}
}{
C_{\max}^{\mathrm{train}}
-
C_{\min}^{\mathrm{train}}
},
\label{eq:vc_normalization}
\end{equation}

with the fitted transformation subsequently applied to the
corresponding evaluation observations.

Two temporal VC-derived features were additionally constructed:
last-observed VC and change from the last available VC observation,

\begin{equation}
VC^{\mathrm{lag}}_{i,t}
=
VC_{i,t-1},
\label{eq:vc_lag}
\end{equation}

\begin{equation}
VC^{\mathrm{change}}_{i,t}
=
VC_{i,t}
-
VC^{\mathrm{lag}}_{i,t}.
\label{eq:vc_change}
\end{equation}

The lag refers to the preceding \emph{available} chronological
observation rather than necessarily the preceding calendar year.
Consequently, for 2020 observations, \(VC^{\mathrm{lag}}\) corresponds
to 2018 because 2019 is absent from the dataset.


\subsubsection{Leakage Control During VC Construction}

All data-dependent quantities used in VC construction, including
rainfall medians, environmental scaling limits, and VC normalization
parameters, were estimated exclusively from the relevant training
partition. The same procedure was repeated within inner training
partitions during model selection. Thus, held-out countries and 2020
evaluation observations did not contribute distributional information
to the construction of their VC-derived features.


\subsection{Chronological Construction of Lagged Features}

Historical malaria-incidence lags were constructed separately within
each country after chronological ordering of the available
observations. Let \(Y_{i,t}\) denote malaria incidence for country
\(i\) at observation \(t\). The most recent and second most recent
available incidence values were defined as

\begin{equation}
L^{(1)}_{i,t}=Y_{i,t-1},
\label{eq:incidence_lag1}
\end{equation}

and

\begin{equation}
L^{(2)}_{i,t}=Y_{i,t-2},
\label{eq:incidence_lag2}
\end{equation}

respectively. These indices refer to the ordering of \emph{available}
observations rather than fixed calendar-year intervals. Thus, because
2019 is absent, \(L^{(1)}_{i,2020}\) corresponds to the observed 2018
incidence.

The one-lag design used \(L^{(1)}\) and retained 127 eligible target
observations, whereas the two-lag design additionally included
\(L^{(2)}\) and retained 95 observations. This comparison allowed the
predictive value of an additional historical observation to be
evaluated against the associated reduction in sample size.

Lag construction used only observations preceding the corresponding
target within each country; the current target incidence was never
used to construct its own historical predictors. For country-grouped
evaluation, a held-out country's previously observed incidence
remained available as a predictor even though its target observations
were excluded from model fitting. The grouped experiment therefore
assesses generalization to held-out countries with historical
epidemiological information available, rather than to countries with
no prior incidence observations.

The same chronological principle was applied to the temporal
VC-derived features defined in
Equations~\ref{eq:vc_lag}--\ref{eq:vc_change}. Accordingly,
\texttt{VC\_lag} represents the last available VC observation and
\texttt{VC\_change} the change from that observation, rather than
necessarily an annual change.



\subsection{Hybrid Feature Representation and Ablation Design}

A structured feature-ablation design was used to quantify the
incremental predictive contribution of different information sources.
Three predictor blocks were distinguished: historical malaria
incidence (\(H\)), environmental and demographic information (\(E\)),
and VC-derived information (\(V\)). Environmental predictors were
represented using either the compact set (\(E_C\)) or the full set
(\(E_F\)), while the VC block comprised

\begin{equation}
V_{i,t}
=
\left[
VC_{\mathrm{norm},i,t},
VC^{\mathrm{lag}}_{i,t},
VC^{\mathrm{change}}_{i,t}
\right].
\label{eq:vc_block}
\end{equation}

Seven configurations were evaluated for each lag design:

\begin{table}[!t]
\centering
\caption{Feature-ablation configurations used in the predictive
experiment.}
\label{tab:ablation_design}
\begin{tabular}{p{0.13\columnwidth}p{0.30\columnwidth}
                p{0.45\columnwidth}}
\toprule
Model & Representation & Information included \\
\midrule
M0 & Persistence & Last-observed incidence prediction \\
M1 & History & Historical incidence only \\
M2 & History + compact environment &
Historical incidence and compact environmental predictors \\
M3 & History + VC &
Historical incidence and VC-derived features \\
M4 & History + compact environment + VC &
Historical incidence, compact environmental predictors, and
VC-derived features \\
M5 & History + full environment &
Historical incidence and full environmental predictors \\
M6 & History + full environment + VC &
Historical incidence, full environmental predictors, and
VC-derived features \\
\bottomrule
\end{tabular}
\end{table}

The incremental contribution of VC was assessed through three direct
ablation contrasts: M3 versus M1 for VC beyond historical incidence,
M4 versus M2 for VC beyond historical incidence and compact
environmental information, and M6 versus M5 for VC beyond historical
incidence and the full environmental representation.

For an error metric \(E\), the corresponding VC contribution was
expressed generally as

\begin{equation}
\Delta_{\mathrm{VC}}
=
E(M_{\mathrm{without\ VC}})
-
E(M_{\mathrm{with\ VC}}),
\label{eq:vc_increment}
\end{equation}

where positive values indicate reduced prediction error after adding
VC-derived information, whereas negative values indicate increased
error.

All non-baseline configurations were evaluated using both Random
Forest and Gradient Boosting, and the same ablation structure was
applied to the one-lag and two-lag datasets. This design isolates the
incremental predictive contribution of VC-derived information from
changes associated with environmental representation or historical
lag structure.

\subsection{Machine-Learning Models}

Two ensemble tree-based regression algorithms were evaluated: Random
Forest (RF) and Gradient Boosting (GB). These models were selected
because they can represent nonlinear relationships and interactions
among historical incidence, environmental, demographic, and VC-derived
predictors without requiring a predefined linear relationship with
malaria incidence.

RF models used 500 trees, a maximum tree depth of 5, a minimum of 3
observations per terminal leaf, and a maximum feature proportion of
0.8. GB models used 180 boosting stages, a learning rate of 0.03, a
maximum tree depth of 2, a minimum of 4 observations per terminal leaf,
and Huber loss. The same model settings were maintained across the
corresponding feature-ablation configurations.

To retain the strong temporal information contained in historical
incidence while allowing additional predictors to contribute, each
machine-learning prediction was blended with the incidence-persistence
prediction:

\begin{equation}
\hat{Y}_{i,t}
=
(1-\alpha)\hat{Y}^{\mathrm{baseline}}_{i,t}
+
\alpha\hat{Y}^{\mathrm{ML}}_{i,t},
\label{eq:prediction_blending}
\end{equation}

where \(\hat{Y}^{\mathrm{baseline}}_{i,t}\) is the persistence
prediction, \(\hat{Y}^{\mathrm{ML}}_{i,t}\) is the RF or GB prediction,
and \(\alpha\) determines their relative contribution. The value of
\(\alpha\) was selected within the training data using
country-grouped inner validation to minimize out-of-sample RMSE.
Neither held-out countries nor 2020 target incidence values were used
in selecting the blending weight.


\subsection{Incidence-Persistence Baseline}

The principal reference model was incidence persistence, for which the
last-observed malaria incidence was used directly as the prediction:

\begin{equation}
\hat{Y}_{i,t}^{\mathrm{baseline}}
=
Y_{i,t-1},
\label{eq:persistence_baseline}
\end{equation}

where \(Y_{i,t-1}\) denotes the most recent available incidence
observation preceding the target for country \(i\). This baseline
provides a strong reference for determining whether environmental and
VC-derived information contributes predictive value beyond the
temporal persistence already present in malaria incidence.


\subsection{Country-Grouped Validation}

Generalization across countries was evaluated using repeated
country-grouped train--test splits. Countries, rather than individual
country-year observations, were used as grouping units so that all
eligible observations from a given country were assigned exclusively
to either the training or evaluation partition within each split. Ten
repeated grouped splits were used.

All data-dependent preprocessing steps, including compact-feature
selection, VC scaling, and selection of the prediction-blending weight,
were fitted using the corresponding training data. Model performance
was then calculated on the held-out countries and summarized across
the repeated splits.

Because historical predictors for a held-out country were constructed
from that country's observations preceding the target, this experiment
evaluates generalization to countries with available historical
incidence information rather than to countries with no epidemiological
history.


\subsection{Temporal Evaluation}

Temporal performance was evaluated by training models on eligible
observations from 2015--2018 and evaluating them on 2020 observations.
The 2020 set contained 31 eligible countries because malaria incidence
for Mali was unavailable in 2020.

All preprocessing operations and model-selection procedures were
restricted to the pre-2020 training data. For 2020 targets, the
last-observed incidence and VC history corresponded to 2018 because
2019 was absent from the dataset.

The 2020 analysis should be interpreted as a temporal holdout rather
than a strict prospective forecast because contemporaneous
environmental information for the target period was included among the
predictors.


\subsection{Performance Metrics}

Predictive performance was assessed using mean absolute error (MAE),
root mean squared error (RMSE), and the coefficient of determination
(\(R^2\)):

\begin{align}
\mathrm{MAE}
&=
\frac{1}{N}
\sum_{j=1}^{N}
\left|y_j-\hat{y}_j\right|,
\\
\mathrm{RMSE}
&=
\sqrt{
\frac{1}{N}
\sum_{j=1}^{N}
\left(y_j-\hat{y}_j\right)^2
},
\\
R^2
&=
1-
\frac{
\sum_{j=1}^{N}(y_j-\hat{y}_j)^2
}{
\sum_{j=1}^{N}(y_j-\bar{y})^2
}.
\end{align}

MAE characterizes average absolute prediction error, RMSE gives greater
weight to larger errors, and \(R^2\) measures the proportion of
variation in incidence explained relative to prediction by the sample
mean. Lower MAE and RMSE and higher \(R^2\) indicate better predictive
performance.


\subsection{Statistical Comparison and Uncertainty}

Uncertainty in temporal model comparisons was assessed using paired
country-level bootstrap resampling. Countries in the 2020 evaluation
set were resampled with replacement, and differences in MAE and RMSE
between paired model configurations were recalculated for each
bootstrap sample. The 2.5th and 97.5th percentiles of the resulting
distributions were used to obtain 95\% bootstrap confidence intervals.

Particular attention was given to the VC ablation contrasts defined
previously (M3 versus M1, M4 versus M2, and M6 versus M5), where
available. For error metrics, positive differences indicate a
reduction in prediction error after adding VC-derived information.

Formal significance tests were not applied to the repeated
country-grouped validation scores because observations and training
sets overlap across repeated splits, violating the independence
assumption underlying conventional paired tests. Grouped-validation
results were therefore summarized descriptively across splits, while
the paired country bootstrap was used to characterize uncertainty in
the temporal evaluation.


\subsection{Feature Importance}

Model-based feature importance was examined as a secondary diagnostic
to characterize the relative contribution of historical,
environmental, and VC-derived predictors. Feature importance was
extracted separately from the fitted Random Forest and Gradient
Boosting models.

Because tree-based importance measures can be influenced by correlated
predictors and do not establish causal effects, they were interpreted
only as descriptive measures of model reliance. The primary evidence
for the predictive contribution of VC was therefore obtained from the
controlled M1--M6 ablation comparisons rather than from feature
importance rankings.


\subsection{Information-Leakage Control and Reproducibility}

All data-dependent transformations were estimated exclusively from the
corresponding training partition. These included environmental
feature selection, rainfall and population scaling used in VC
construction, VC normalization, and selection of the prediction
blending weight. During inner validation, these operations were
re-estimated using only the relevant inner training data before being
applied to validation observations.

Lagged incidence and VC features were constructed chronologically
within countries using only information available before the
corresponding target observation. Held-out target incidence values and
2020 target incidence values were therefore never used in feature
construction, preprocessing, or model selection.

The same preprocessing, feature-ablation definitions, model settings,
and evaluation procedures were applied consistently to Random Forest
and Gradient Boosting models. Random seeds were fixed for repeated
data splitting and model fitting to support computational
reproducibility.


\section{Results}

\subsection{Dataset Characteristics}

The target-complete annual dataset contained 159 country-year
observations from 32 Sub-Saharan African countries. The one-lag design
retained 127 eligible observations from 2016, 2017, 2018, and 2020,
whereas the two-lag design retained 95 observations from 2017, 2018,
and 2020. Because 2019 was absent, lagged information for 2020
corresponded to the most recent available observation in 2018.


\subsection{Grouped Validation Performance}

Repeated country-grouped validation showed generally lower prediction
errors for the one-lag than the two-lag design. Table~\ref{tab:grouped_results}
summarizes the principal one-lag Random Forest results. Persistence
(M0) achieved the lowest mean MAE (22.2711), whereas the
history-plus-compact-environment model (M2) achieved the lowest RMSE
(37.3853) and highest \(R^2\) (0.8480). Thus, the relative performance
of the models depended on the evaluation metric.

\begin{table}[!t]
\centering
\caption{Repeated country-grouped validation performance for the
one-lag Random Forest configurations.}
\label{tab:grouped_results}
\small
\setlength{\tabcolsep}{4pt}
\begin{tabular}{lccc}
\toprule
Model & MAE & RMSE & $R^2$ \\
\midrule
Persistence (M0)      & \textbf{22.2711} & 39.8300 & 0.8268 \\
History RF (M1)       & 23.5992 & 39.0749 & 0.8336 \\
Hist.+Env. RF (M2)    & 24.1465 & \textbf{37.3853} & \textbf{0.8480} \\
Hist.+VC RF (M3)      & 24.2671 & 38.1615 & 0.8429 \\
Hist.+Env.+VC RF (M4) & 24.2333 & 37.4594 & 0.8477 \\
Full Env. RF (M5)     & 24.4522 & 37.8454 & 0.8443 \\
Full Env.+VC RF (M6)  & 24.3378 & 37.8552 & 0.8443 \\
\bottomrule
\end{tabular}
\end{table}

Among the one-lag Random Forest configurations, M2 achieved the lowest
grouped RMSE (37.3853) and highest \(R^2\) (0.8480), while M4 produced
similar performance (RMSE 37.4594; \(R^2=0.8477\)). The more
parsimonious M3 configuration achieved an RMSE of 38.1615 and
\(R^2=0.8429\). In contrast, persistence retained the lowest MAE
(22.2711), indicating that the machine-learning gains in RMSE were
concentrated more strongly among larger prediction errors.

The corresponding two-lag M3 Random Forest model produced a grouped
RMSE of 42.1408, compared with 38.1615 for the one-lag model. Thus,
the additional historical lag did not compensate for the reduction
from 127 to 95 eligible observations.


\subsection{Temporal Evaluation on 2020 Observations}

Temporal evaluation included 31 countries with available 2020 malaria
incidence. Table~\ref{tab:temporal_results} summarizes the principal
one-lag Random Forest results.

\begin{table}[!t]
\centering
\caption{2020 temporal-evaluation performance for the one-lag Random
Forest configurations.}
\label{tab:temporal_results}
\small
\setlength{\tabcolsep}{4pt}
\begin{tabular}{lccc}
\toprule
Model & MAE & RMSE & $R^2$ \\
\midrule
Persistence (M0)      & 32.2710 & 54.6111 & 0.6571 \\
History RF (M1)       & 31.5723 & 52.7665 & 0.6798 \\
Hist.+Env. RF (M2)    & 31.3547 & 52.4761 & 0.6834 \\
Hist.+VC RF (M3)      & \textbf{30.5827} & \textbf{50.4361} & \textbf{0.7075} \\
Hist.+Env.+VC RF (M4) & 31.5876 & 52.3498 & 0.6849 \\
Full Env. RF (M5)     & 31.5203 & 52.7202 & 0.6804 \\
Full Env.+VC RF (M6)  & 31.5313 & 52.5652 & 0.6823 \\
\bottomrule
\end{tabular}
\end{table}

The one-lag history-plus-VC Random Forest model (M3) achieved the
lowest MAE and RMSE and the highest \(R^2\) among the evaluated
configurations, with an MAE of 30.5827, RMSE of 50.4361, and
\(R^2\) of 0.7075. Relative to persistence (M0), M3 reduced MAE by
approximately 5.2\% and RMSE by 7.6\%, while increasing \(R^2\) from
0.6571 to 0.7075.

For the direct VC ablation, M3 reduced RMSE by 2.3305 incidence points
and MAE by 0.9896 relative to the corresponding history-only Random
Forest model (M1). In paired country-level bootstrap analysis, the
RMSE reduction had a 95\% interval of $[-0.4516,\,4.0317]$, with
95.45\% of bootstrap resamples favoring M3 over M1; the corresponding
MAE interval was $[-0.7553,\,2.7919]$. Although both intervals included
zero, these results provide the clearest point-estimate evidence of
incremental predictive information from the VC-derived representation
beyond historical incidence alone.

The incremental differences were more modest when environmental
predictors were already included. RMSE decreased from 52.4761 in M2
to 52.3498 in M4 and from 52.7202 in M5 to 52.5652 in M6. Thus, the
largest point-estimate reduction in temporal RMSE associated with
adding VC occurred in the history-only comparison (M3 versus M1).

The two-lag M3 model produced an RMSE of 50.8512 and \(R^2=0.7027\),
compared with 50.4361 and 0.7075 for the one-lag formulation.
Together with the larger analytical sample retained by the one-lag
design, these results supported its use as the primary analysis.

\subsection{Prediction Diagnostics}

Prediction diagnostics for the 2020 temporal holdout are shown in
Figures~\ref{fig:observed_predicted}--\ref{fig:error_boxplot}.
Figure~\ref{fig:observed_predicted} compares observed incidence with
predictions from the persistence baseline and the one-lag
history-plus-VC Random Forest model (M3). The M3 predictions showed
lower overall error than persistence, although larger discrepancies
remained for some high-incidence observations.

\begin{figure}[!tbp]
\centering
\includegraphics[width=\columnwidth]
{Figure_2_observed_vs_predicted_corrected.png}
\caption{Observed versus predicted malaria incidence in the 2020
temporal holdout for the incidence-persistence baseline and the
one-lag history-plus-vectorial-capacity Random Forest model (M3).
The dashed diagonal represents perfect agreement.}
\label{fig:observed_predicted}
\end{figure}

Figure~\ref{fig:error_distribution} shows the cumulative distribution
of absolute prediction errors. At an error threshold of 25 incidence
points, 58.1\% of countries were within the threshold for M3,
compared with 54.8\% for both persistence and the corresponding
history-plus-VC Gradient Boosting model.

\begin{figure}[!tbp]
\centering
\includegraphics[width=\columnwidth]
{Figure_3_error_distributions_corrected.png}
\caption{Cumulative distribution of absolute prediction errors in the
2020 temporal holdout. Curves rising more rapidly indicate a larger
proportion of countries with smaller absolute errors.}
\label{fig:error_distribution}
\end{figure}

Figure~\ref{fig:error_boxplot} further illustrates the distribution
and variability of prediction errors across models, showing that
aggregate performance differences were not associated with uniform
improvement across all countries.

\begin{figure}[!tbp]
\centering
\includegraphics[width=\columnwidth]
{Figure_4_error_boxplot_corrected.png}
\caption{Distribution of prediction errors across the evaluated models
in the 2020 temporal holdout.}
\label{fig:error_boxplot}
\end{figure}

\subsection{Incremental Contribution of Vectorial Capacity}

The ablation analysis showed that the predictive contribution of
VC-derived information depended on the accompanying feature
representation. In the 2020 one-lag Random Forest evaluation, adding
VC-derived features directly to historical incidence (M3 versus M1)
reduced RMSE from 52.7665 to 50.4361 and MAE from 31.5723 to 30.5827,
corresponding to reductions of 2.3305 and 0.9896 incidence points,
respectively.

When compact environmental predictors were already included (M4
versus M2), the RMSE reduction was 0.1262 incidence points
(95\% bootstrap CI: $-0.3626$ to $0.3651$), with 72.35\% of bootstrap
resamples favoring M4. The corresponding MAE reduction was $-0.2329$
(95\% CI: $-0.6228$ to $0.1639$).

Similarly, adding VC to the full environmental representation (M6
versus M5) reduced RMSE by 0.1551 incidence points
(95\% CI: $-0.1508$ to $0.3262$), with 87.95\% of bootstrap resamples
favoring M6. The corresponding MAE reduction was $-0.0110$
(95\% CI: $-0.3334$ to $0.2766$).

Gradient Boosting showed similarly small incremental RMSE differences
when VC was added to environmental representations: 0.3592
(95\% CI: $-0.2073$ to $0.7366$) for the compact representation and
0.1271 (95\% CI: $-0.1229$ to $0.3531$) for the full representation.

Thus, the clearest point-estimate improvement from VC-derived
information occurred when it was combined directly with historical
incidence in the Random Forest model. When compact or full
environmental predictors were already present, the incremental
differences were small and the bootstrap intervals included zero.

\subsection{Country-Level Heterogeneity}

Country-level performance under repeated grouped evaluation was
heterogeneous. The one-lag history-plus-VC Random Forest model reduced
mean absolute error relative to persistence in 7 of the 29 countries
with paired comparisons (Figure~\ref{fig:country_rf}). Thus, the
aggregate grouped-validation performance did not correspond to uniform
country-level improvements.

\begin{figure}[!t]
\centering
\includegraphics[width=\columnwidth]
{Figure_7_country_improvement_hybrid_rf.png}
\caption{Country-level mean absolute-error improvement for the one-lag
history-plus-vectorial-capacity Random Forest model relative to
persistence under repeated country-grouped evaluation. Positive values
indicate lower error for the history-plus-VC model.}
\label{fig:country_rf}
\end{figure}

The corresponding Gradient Boosting model reduced mean absolute error
relative to persistence in 8 of 29 countries
(Figure~\ref{fig:country_gb}), again indicating that predictive gains
were concentrated in a subset of countries.

\begin{figure}[!t]
\centering
\includegraphics[width=\columnwidth]
{Figure_7_country_improvement_hybrid_gb.png}
\caption{Country-level mean absolute-error improvement for the one-lag
history-plus-vectorial-capacity Gradient Boosting model relative to
persistence under repeated country-grouped evaluation. Positive values
indicate lower error for the history-plus-VC model.}
\label{fig:country_gb}
\end{figure}


\subsection{Importance of VC-Derived Features}

Feature importance provided a complementary view of how VC-derived
information was used in the final one-lag history-plus-VC models.
As shown in Figure~\ref{fig:vc_importance}, \texttt{VC\_norm} was the
most important VC-derived feature in both algorithms. Its importance
was 0.0805 in Random Forest and 0.0263 in Gradient Boosting, compared
with 0.0431 and 0.0055 for \texttt{VC\_lag}, and 0.0266 and 0.0092
for \texttt{VC\_change}, respectively.

Historical incidence remained the dominant predictor, with
\texttt{lag1\_incidence} importance of 0.8498 in Random Forest and
0.9591 in Gradient Boosting. Thus, although VC-derived features
contributed to the fitted models, their individual importance was
substantially smaller than that of recent malaria incidence. These
importance values are interpreted as secondary diagnostics; the
ablation experiments provide the primary evidence for the incremental
predictive contribution of VC.

\begin{figure}[!tbp]
\centering
\includegraphics[width=\columnwidth]{Figure_9_vc_feature_importance.png}
\caption{Impurity-based importance of VC-derived features in the final
one-lag history-plus-VC Random Forest and Gradient Boosting models for
the 2020 temporal evaluation. Values indicate predictor utilization
within the fitted models and not causal influence.}
\label{fig:vc_importance}
\end{figure}


\section{Discussion}

\subsection{Principal Findings}

This study evaluated whether vectorial-capacity-derived information
provides incremental predictive value when integrated into a
machine-learning framework for malaria-incidence prediction. The
results indicate a context-dependent contribution that varied with
feature representation, learning algorithm, evaluation strategy, and
performance metric.

The clearest improvement occurred in the 2020 temporal evaluation.
The one-lag history-plus-VC Random Forest model (M3) achieved an RMSE
of 50.4361 compared with 54.6111 for persistence, corresponding to a
7.6\% reduction, while \(R^2\) increased from 0.6571 to 0.7075.
More importantly, relative to the corresponding history-only Random
Forest model (M1), inclusion of VC-derived information reduced RMSE by
2.3305 and MAE by 0.9896 incidence points. This indicates incremental
predictive information beyond historical incidence alone in this
configuration.

The contribution was substantially smaller when compact or full
environmental predictors were already included, and the bootstrap
intervals for these VC-ablation contrasts included zero. This pattern
is consistent with partial information overlap because the VC
representation is itself constructed from environmental and
demographic variables.

Country-grouped evaluation provided a complementary result: several
machine-learning configurations reduced RMSE relative to persistence,
whereas persistence retained the lowest mean MAE. The one-lag design
also generally outperformed the two-lag design while retaining 127
rather than 95 eligible observations. Overall, the findings support a
conditional role for VC-derived information, with its clearest value
in the present analysis occurring as a compact biologically informed
representation combined with recent malaria incidence.

\subsection{Biological Interpretation and Incremental Value of VC}

The contribution of this study is not a new vectorial-capacity
formulation, but the integration and evaluation of environmentally
parameterized VC within a machine-learning framework. Classical
vectorial capacity organizes transmission potential through
components related to mosquito abundance, biting activity, survival,
and parasite development
\parencite{Ross1905,Ross1907,LindsayBirley1996}. In the present study,
these components were represented using temperature, rainfall,
vegetation, and population information rather than direct
entomological measurements.

This distinction is central to interpretation. The resulting
VC-derived features should not be regarded as direct country-specific
estimates of mosquito density, biting rate, survival probability, or
parasite-development duration. Instead, they provide a biologically
structured representation of environmental conditions associated with
transmission potential. Temperature influences several mosquito and
parasite processes, while rainfall and vegetation contribute to
conditions relevant to vector populations
\parencite{Villena2022,Dahlin2024}.

The ablation results suggest that this biological representation can
complement historical epidemiological information. Historical
incidence summarizes the recent observed state of malaria transmission
and implicitly captures multiple measured and unmeasured determinants,
whereas VC organizes selected environmental and demographic conditions
through an epidemiologically motivated structure. The improved
temporal performance of M3 relative to M1 indicates that these
information sources were not completely redundant in the evaluated
Random Forest configuration.

The smaller and uncertain improvements obtained when VC was added to
compact or full environmental representations are also informative.
Because VC is constructed partly from the same environmental and
demographic predictors, broader environmental models may already
contain much of the statistical information summarized by the VC
representation. Its contribution may therefore lie not only in adding
new predictive information, but also in organizing existing
environmental information into a compact, biologically interpretable
representation. This is consistent with the broader objective of
knowledge-guided and mechanism-informed machine learning, in which
domain knowledge constrains or structures data-driven prediction
\parencite{Yin2019,Wang2020,Karpatne2024}.

The observed predictive benefit should not, however, be interpreted as
validation of the calculated VC values as direct measures of biological
transmission intensity. Such validation would require independent
entomological observations of vector abundance, biting behavior,
survival, and related transmission components. The present findings
instead demonstrate the predictive utility of an environmentally
parameterized vectorial-capacity representation within the evaluated
machine-learning framework.


\subsection{Historical Incidence and the Contribution of VC-Derived Features}

Historical incidence remained the dominant predictor in the final
one-lag history-plus-VC models. The importance of
\texttt{lag1\_incidence} was 0.8498 in Random Forest and 0.9591 in
Gradient Boosting, substantially exceeding the importance of the
individual VC-derived features. This dominance is consistent with the
strong temporal persistence of malaria incidence, which can summarize
the combined influence of previous transmission intensity,
interventions, healthcare access, surveillance characteristics, and
other country-specific factors not explicitly represented in the
predictor set.

Among the VC-derived features, \texttt{VC\_norm} had the highest
importance in both algorithms. Its importance was 0.0805 in Random
Forest and 0.0263 in Gradient Boosting, compared with 0.0431 and
0.0055 for \texttt{VC\_lag}, and 0.0266 and 0.0092 for
\texttt{VC\_change}, respectively. This pattern indicates that the
models relied more strongly on the contemporaneous normalized VC
representation than on its historical level or change from the last
available observation.

The dominance of historical incidence also explains the strength of
the persistence baseline, particularly under country-grouped
evaluation where it retained the lowest mean MAE. Nevertheless,
historical dominance does not imply that VC-derived information was
uninformative. The temporal ablation results showed that adding VC to
historical incidence improved the one-lag Random Forest prediction,
supporting a complementary rather than competing relationship between
the two information sources.

Feature importance should be interpreted cautiously because correlated
predictors can redistribute importance and tree-based importance does
not establish causal influence. Accordingly, these results are used
only to describe predictor utilization within the fitted models; the
controlled feature-ablation experiments remain the primary evidence
for the incremental predictive contribution of VC-derived information.


\subsection{Generalization and Country-Level Heterogeneity}

The country-grouped and temporal evaluations examined complementary
dimensions of predictive generalization. Country-grouped validation
assessed performance in held-out countries with historical incidence
information available, whereas the 2020 temporal holdout assessed
performance during a later observation period using models developed
from earlier data.

The two evaluation settings produced different performance patterns.
Under country-grouped validation, several machine-learning
configurations reduced RMSE and increased \(R^2\) relative to
persistence, with the history-plus-compact-environment Random Forest
model (M2) achieving the lowest grouped RMSE. Persistence nevertheless
retained the lowest mean MAE. In the temporal holdout, the
history-plus-VC Random Forest model (M3) achieved the lowest RMSE,
together with lower MAE and higher \(R^2\) than persistence. These
differences demonstrate that conclusions about predictive performance
depend on both the validation structure and the performance metric.

Country-level analysis further showed that predictive gains were not
geographically uniform. Under repeated country-grouped evaluation,
the history-plus-VC Random Forest model reduced mean absolute error
relative to persistence in 7 of 29 countries, while the corresponding
Gradient Boosting model did so in 8 of 29 countries. This
heterogeneity is consistent with variation in malaria epidemiology,
environmental conditions, interventions, healthcare and surveillance
systems, and historical transmission patterns across countries
\parencite{Odhiambo2020,Wu2023}.

These findings emphasize the importance of evaluating malaria
prediction models across both geographical and temporal dimensions
rather than relying on a single aggregate validation result. They also
indicate that the predictive contribution of VC-derived information
should not be assumed to generalize uniformly across epidemiological
settings.

The 2020 experiment should nevertheless be interpreted as a temporal
holdout rather than a strictly prospective forecast because
contemporaneous environmental predictors from 2020 were included in
the predictor representation. Historical incidence and lagged VC
information remained restricted to observations preceding the target
period.


\subsection{Relationship to Previous Work and Implications}

Previous malaria-prediction studies have demonstrated the value of
climatic, environmental, remote-sensing, and epidemiological predictors,
including temperature, rainfall, vegetation, and historical incidence
\parencite{MafweleLee2022,Okiring2021,Cuellar2025,Odu2021,Masinde2020,Kaur2022}.
The present study extends this direction by incorporating
vectorial-capacity-derived information as a biologically structured
representation within the predictive feature space rather than simply
adding further environmental variables.

The proposed framework is therefore best characterized as biologically
informed machine learning rather than as a new mechanistic model of
malaria transmission. The classical VC formulation provides
epidemiological structure, while Random Forest and Gradient Boosting
learn predictive relationships among historical incidence,
environmental conditions, VC-derived features, and observed malaria
incidence. This approach is consistent with the broader principle of
integrating domain knowledge into data-driven models without requiring
the complete prediction process to be specified mechanistically
\parencite{Jhutty2022,Karpatne2024}.

The results also demonstrate that biological structure does not
guarantee improved prediction in every configuration. VC-derived
information provided its clearest contribution when combined directly
with recent incidence, whereas its incremental benefit was smaller and
uncertain when broader environmental representations were already
included. This suggests that biologically informed representations may
be particularly useful as compact complements to strong epidemiological
predictors rather than as universally additive feature sets.

For malaria surveillance and predictive decision support, this
framework provides a proof-of-concept for combining recent
epidemiological information with biologically structured environmental
information. However, the observed country-level heterogeneity
indicates that performance should be evaluated locally before
operational application. Future work should assess the framework using
larger and more continuous datasets, finer spatial and temporal
resolution, independent geographical validation, and direct
entomological measurements where available. These extensions would
permit stronger evaluation of the relationship between environmentally
parameterized vectorial capacity and observed malaria transmission and
support the development of more localized prediction and early-warning
systems.



\section{Limitations and Future Research}

\subsection{Data Scale, Resolution, and Generalizability}

The present study was based on 159 target-complete annual
country-level observations from 32 Sub-Saharan African countries, with
127 observations eligible for the primary one-lag analysis and 95 for
the two-lag analysis. This relatively small sample constrained model
complexity, while annual country-level aggregation masked
within-country spatial and seasonal variation in environmental
conditions, vector ecology, interventions, and malaria transmission.
Furthermore, because 2019 was absent, the last-observed information
used for 2020 corresponded to 2018 rather than the immediately
preceding calendar year.

These limitations motivate evaluation using larger and more continuous
surveillance datasets at finer spatial and temporal resolutions.
External validation across independent countries, ecological zones,
transmission settings, and time periods will also be important for
determining the generalizability of the observed predictive
relationships.

The 2020 experiment should additionally be interpreted as a temporal
holdout rather than a strictly prospective forecast because
contemporaneous environmental predictors were included. Future
forecasting applications should restrict predictors to information
available before the prediction period or incorporate forecast and
near-real-time environmental data alongside epidemiological
surveillance.

\subsection{VC Representation and Future Model Development}

The vectorial-capacity representation was environmentally parameterized
from temperature, rainfall, vegetation, and population information
rather than from direct country-specific entomological measurements.
Accordingly, the resulting VC-derived features represent biologically
structured transmission-related environmental information and should
not be interpreted as direct measurements of mosquito abundance,
biting behavior, survival, parasite development, or vectorial capacity
itself.

Future studies should incorporate direct entomological observations,
including vector abundance and species composition, biting rates,
survival, and parasite-development measurements, where available.
Combined with higher-resolution environmental and epidemiological data,
such observations would permit more empirically grounded
parameterization and validation of VC components and stronger
assessment of their predictive contribution across transmission
settings.

Random Forest and Gradient Boosting were selected because they provide
flexible nonlinear prediction while remaining suitable for the limited
analytical sample. Larger datasets would permit evaluation of
additional machine-learning and hybrid approaches while preserving the
biologically informed structure of the framework. Such extensions,
together with independent external validation, could support future
development toward localized malaria prediction and operational
early-warning systems.


\section{Conclusion}

This study evaluated the incremental predictive contribution of
vectorial-capacity-derived information within a biologically informed
machine-learning framework for malaria-incidence prediction. Across
the evaluated configurations, the contribution of VC was
context-dependent rather than universal. The clearest evidence was
observed in the 2020 temporal evaluation, where the one-lag
history-plus-VC Random Forest model achieved the lowest temporal RMSE
(50.4361) among the evaluated configurations. Relative to the
corresponding history-only model, incorporating VC-derived information
reduced RMSE by 2.33 incidence points and MAE by 0.99 incidence
points, demonstrating incremental predictive value beyond historical
incidence alone in this setting. The broader ablation analysis showed that the magnitude of the incremental VC contribution varied across feature representations,
while country-grouped performance also differed across metrics and
countries. Historical incidence remained the dominant predictive
signal, with \texttt{VC\_norm} providing the strongest contribution
among the VC-derived features. These findings suggest that the
principal value of VC in the present framework lies in providing a
compact, biologically structured representation that can complement
recent epidemiological information, particularly when extensive
environmental predictors are not already represented in the model.

The study therefore provides a proof-of-concept for integrating
epidemiological knowledge into data-driven malaria prediction without
treating the resulting model as a mechanistic simulation or the
environmentally parameterized VC representation as direct
entomological measurement. Further evaluation using larger and more
continuous datasets, finer spatio-temporal resolution, independent
external validation, and direct entomological observations is needed
to establish the generalizability and potential operational value of
this biologically informed approach.

\printbibliography[title={References}]

\end{document}